\BookPart{Reliability, scale and storage}{reliability}{%
Molecules are noisy: strands mis-pair, reactions leak, assemblies grow the
wrong tile, and readouts blur. Part~VI treats error as a first-class design
quantity. It develops fault models and redundancy, asks which resources grow
when a computation is scaled, and then turns to the most commercially advanced
use of DNA as an information medium: archival data storage, where
error-correcting codes, synthesis and sequencing meet. It closes with the
laboratory practices that make a molecular result reproducible.

\textbf{Chapters in this part}
\begin{itemize}\setlength{\itemsep}{2pt}
\item \textbf{23\quad Noise, crosstalk and molecular fault models.} How do local errors accumulate, and how can redundancy and proofreading suppress them?
\item \textbf{24\quad Scaling and resource limits.} Which costs grow with problem size, and why does parallelism not make them vanish?
\item \textbf{25\quad DNA storage and error-correcting design.} How are files written into, retrieved from and recovered out of a pool of synthetic DNA?
\item \textbf{26\quad Laboratory workflows and reproducibility.} What must be recorded for a molecular computation to count as evidence?
\end{itemize}}
